\documentclass[a4paper,twoside]{article}


% --- RUIMTE EN LAYOUT ---
\usepackage{setspace}

\usepackage{geometry} %size regulating
\geometry{
 a4paper,
 left=25mm,
 right=25mm,
 top=20mm,
 bottom=20mm,
 }

\usepackage{parskip} % for better spacing between paragraphs
\usepackage{flushend} % Alleen nodig als je laatste pagina van artikel wilt vullen met tekst, zodat de kolommen even lang zijn. Anders kan je dit beter weglaten.

% --- WISKUNDE ENSYMBOLEN ---
\usepackage{amsmath, amssymb, amsfonts, amsthm} % For nicer math format
\usepackage{mathrsfs} % For curly arrr
\usepackage{bm} % For bold symbols eg \bm{\nabla}
\usepackage{esint} % For bolletje in integrals
\usepackage{dsfont}
\usepackage{siunitx} % Voeg deze expliciet toe voor die graden en eenheden!
\usepackage{braket} % For bra-ket notation
\usepackage{mathtools}
\usepackage{url}
\allowdisplaybreaks

% --- OVERIGE ---

\usepackage{graphicx} % Required for inserting images
\usepackage[hidelinks]{hyperref} % geen kaders om linken
\usepackage{subcaption}
%\usepackage[dutch]{babel}
\usepackage[T1]{fontenc}
\usepackage[utf8]{inputenc}
\usepackage{lmodern}
\usepackage{multicol}


% --- Tikz packages ---
\usepackage{tikz}
\usetikzlibrary{angles,quotes, babel} % for pic (angle labels)
\usetikzlibrary{arrows.meta,decorations.markings,calc, decorations.pathmorphing, decorations.pathreplacing, positioning, shapes, shadings, patterns}
\usepackage{xcolor}
\usepackage{tikz-3dplot}
\usepackage[siunitx]{circuitikz}

\usepackage{etoolbox} % ifthen
\usepackage[outline]{contour} % glow around text
\usepackage{xfp} % higher precision (16 digits?)
\contourlength{1.1pt}



% --- THEOREM STIJLEN MET EXTRA RUIMTE ---

\theoremstyle{definition}
\newtheorem{postulate}{Postulaat}
\newtheorem{definition}{Definitie}
\newtheorem{derivation}{Afleiding}
\newtheorem{proposition}{Propositie}
\newtheorem{example}{Voorbeeld}

\theoremstyle{plain}
\newtheorem{theorem}{Theorem}


% --- EXTRA COMMANDO'S ---

\newcommand{\dbar}{{\mkern+3mu\mathchar'26\mkern-12mu d}}

% --- Titel en Introductie --- 

\title{Statistische Fysica - Mijn oefeningen}
\author{T. Cools}
\date{2026-2027}

\begin{document}
\maketitle
\begin{abstract}
    Dit document bevat mijn persoonlijke oefeningen voor het vak Statistische Fysica.
\end{abstract}

\section{Reeks 1}
\subsection{Oefening 1}
Deze oefening is herhaling van thermische fysica.
\subsection{Oefening 2}
\subsubsection{deel 1}
Grenzen zijn correct!
\subsubsection{deel 2}
\[
S = k\left[N\ln(2\mu) - \frac{3 M^2}{2N\mu^2}\right]
\]
Entropie juist!
\subsubsection{deel 3}
\[
M(B,T) = \left(\frac{N \mu^2}{3k}\right) \frac{B}{T}
\]
was correct.
\subsubsection{deel 4}
Vind ik wat vaag...
\subsection{Oefening 3}
\subsubsection{Deel 1}
Entropie is juist.
\subsubsection{Deel 2}
correct
\subsubsection{Deel 3}
Crazy berekeningen die ik ook correct heb!
\subsubsection{Deel 4}
Crazy berekingen waar ik jammer genoeg een foutje heb gemaakt.
\subsubsection{Deel 5}
Niet meer juist
\subsection{Oefening 4}
Ik had de referentiie energie niet verelegd naar het nulpunt, waardoor mijn berekeningen vaan complexer zijn.
Daardoor heb ik een fout gemaakt, ik zal deze oefening nog eens opnieuw moeten maken.
\subsection{Oefening 5}
\subsubsection{Deel 1}
In principe wel correct, meer ik heb de berekening niet volledig uitgevoerd:
\[
\sum_{r} e^{-0.3\beta r  } = \frac{1}{1 - e^{-0.3\beta}}
\]
\subsection{Oefening 6}
\subsubsection{Deel 1}
Weer een klein rekenfoutje...
\subsubsection{Deel 2}
Nu waren de (moeilijkere) berekeningen dan wel correct.
\subsection{Oefening 7}
Fout, we moeten de voorwaarde vinden zodat \(S\) inderdaad een extensieve grootheid is.
\[
S(\lambda E , \lambda L, \lambda V) = \lambda S(E, L, V)
\]
\section{Reeks 2}
\end{document}